% ============================================================
% РОЗДІЛ 2 — переписаний (структура без підпідрозділів, стиль академічний)
% Вставити замість блоку від \protect\phantomsection\label{_Toc228121443}{}
% до \protect\phantomsection\label{_Toc229681122}{} у БКР_Тесля.tex
% ============================================================

\protect\phantomsection\label{_Toc228121443}{}РОЗДІЛ 2. ПРОЄКТУВАННЯ
ПРОГРАМНОГО РІШЕННЯ

\protect\phantomsection\label{_Toc228121444}{}2.1. Вибір методів,
технологій та засобів реалізації

У підрозд.~1.5 сформульовано технічне завдання: розробити веб-симулятор
потокового обчислювача 16-точкового ШПФ з інтерактивною візуалізацією
руху токенів, числовою верифікацією та відкритим доступом. У цьому
підрозділі обґрунтовується вибір архітектурних підходів, парадигм
реалізації та сімейств інструментів, що покривають зазначені вимоги.
Конкретні версії бібліотек та параметри середовища розробки наводяться
у підрозд.~3.1.

Для визначення основного архітектурного підходу розглянуто три
альтернативи. \textbf{Варіант~А --- настільний застосунок} (Electron,
Qt, JavaFX) забезпечує повний доступ до ресурсів хост-системи та
багатопоточність, однак вимагає встановлення на кожну робочу станцію,
окремих збірок під різні ОС і суперечить вимозі NF7 (робота без
встановлення). \textbf{Варіант~Б --- клієнт-серверний веб-застосунок}
передбачає розміщення ядра симуляції на сервері, що вносить мережеву
затримку, несумісну з вимогою NF2 (час відгуку \(\leq 100\)~мс), і
потребує підтримки серверної інфраструктури. \textbf{Варіант~В ---
повністю клієнтський SPA (Single-Page Application)} обрано як основу
архітектури: усі обчислення виконуються у браузері, серверна частина
відсутня, застосунок завантажується одноразово як статичний пакет
(HTML+JS+CSS) і надалі не потребує мережі. Цей варіант одночасно
задовольняє вимоги NF2, NF4 (портативність), NF5 (детермінованість) і
NF7.

Симулятор поєднує три принципово різні підсистеми: анімована графова
візуалізація DFG, числове ядро симуляції та панель керування і метрик.
Для їх ефективної координації обрано поєднання \emph{реактивного
рендерингу} та \emph{централізованого спостерігаваного сховища стану}.
Реактивний рендеринг (React-парадигма) природно відображає стан
симулятора на візуальні компоненти: зміна значення в одному полі сховища
автоматично ініціює перемалювання лише тих компонентів, що підписані на
це поле, що забезпечує стабільні 60~fps (NF3). Централізоване сховище
стану (pattern single source of truth) гарантує, що всі панелі ---
GraphView, SpectrumView, Metrics --- бачать один і той самий «знімок»
стану у поточному кадрі. Альтернатива у вигляді прямого маніпулювання
DOM (jQuery-style) відкинута через лінійне зростання складності, а
Flux-підходи з ред'юсерами (Redux) --- через надмірний обсяг
шаблонного коду.

Для реалізації інтерактивного графа потоку даних розглянуто рендеринг
на HTML-canvas (максимальна продуктивність, але ручна реалізація всіх
UI-взаємодій), SVG з ручним кодуванням (простіший налагодження, але
обмежена розширюваність) та спеціалізовану бібліотеку графової
візуалізації React~Flow. Обрано React~Flow: бібліотека надає готові
механізми перетягування вузлів, зума, анімованих ребер і плаваючих
підказок, а формальна модель \texttt{Node}/\texttt{Edge} добре
відображається на абстракцію DFG.

У проєктуванні системи застосовано п'ять класичних шаблонів
проєктування. \emph{Observer (Publisher/Subscriber)} забезпечує розсилку
подій \texttt{onToken}, \texttt{onFire}, \texttt{onOutput}: ядро
симуляції є видавцем, а UI-компоненти --- підписниками без прямих
посилань на ядро. \emph{State~/ Finite State Machine} керує життєвим
циклом симуляції:
\texttt{IDLE}\(\rightarrow\)\texttt{RUNNING}\(\rightarrow\)\texttt{PAUSED}\(\rightarrow\)\texttt{DONE}.
\emph{Strategy} забезпечує різні типи вузлів (\texttt{source},
\texttt{dft4}, \texttt{twiddle}, \texttt{sink}): поведінка вузла
інкапсулюється у функції, обраній за полем \texttt{NodeKind}. \emph{Factory}
реалізовано через функцію \texttt{generateDFT4x4()}, що повертає готовий
об'єкт \texttt{Graph} за параметрами \(N\) та \texttt{latency}.
\emph{Facade} --- клас \texttt{DataflowRuntime} приховує внутрішню
реалізацію черг і планувальника за єдиним публічним методом
\texttt{tick(dt,~maxFires)}.

Для дотримання NF1 (точність \(\varepsilon \leq 10^{-9}\)) і NF5
(детермінованість) передбачено статичну типізацію вихідного коду, що
унеможливлює помилки передачі аргументів неправильного типу на етапі
компіляції. Комплексні числа подаються структурою \(\{re,\,im\}\) у
подвійній точності IEEE~754 (\(\varepsilon_{mach} \approx 2.22 \cdot 10^{-16}\)),
що забезпечує необхідний запас точності.

Підсумкова таблиця зв'язків між вимогами підрозд.~1.5 та обраними
архітектурними рішеннями наведена у табл.~2.1.

Таблиця 2.1.

Зв'язок вимог і архітектурних рішень

\begin{longtable}[]{@{}
  >{\raggedright\arraybackslash}p{(\linewidth - 4\tabcolsep) * \real{0.2518}}
  >{\raggedright\arraybackslash}p{(\linewidth - 4\tabcolsep) * \real{0.3630}}
  >{\raggedright\arraybackslash}p{(\linewidth - 4\tabcolsep) * \real{0.3852}}@{}}
\caption[Зв'язок вимог і архітектурних рішень]{Зв'язок вимог і
архітектурних рішень}\tabularnewline
\toprule\noalign{}
\begin{minipage}[b]{\linewidth}\raggedright
\textbf{Вимога}
\end{minipage} & \begin{minipage}[b]{\linewidth}\raggedright
\textbf{Архітектурне рішення}
\end{minipage} & \begin{minipage}[b]{\linewidth}\raggedright
\textbf{Обґрунтування}
\end{minipage} \\
\midrule\noalign{}
\endfirsthead
\toprule\noalign{}
\begin{minipage}[b]{\linewidth}\raggedright
\textbf{Вимога}
\end{minipage} & \begin{minipage}[b]{\linewidth}\raggedright
\textbf{Архітектурне рішення}
\end{minipage} & \begin{minipage}[b]{\linewidth}\raggedright
\textbf{Обґрунтування}
\end{minipage} \\
\midrule\noalign{}
\endhead
\bottomrule\noalign{}
\endlastfoot
NF1 (точність) & Статична типізація + IEEE~754 double & Виключення
помилок типів, \(\varepsilon_{mach} \approx 2.22 \cdot 10^{-16}\) \\
NF2 (відгук \(\leq 100\)~мс) & Клієнтський SPA без серверних викликів &
Відсутність мережевої затримки \\
NF3 (60~fps) & Реактивний рендеринг з підпискою на окремі поля &
Перемальовування лише змінених компонентів \\
NF4 (портативність) & Стандарт ES2020, Web APIs & Робота у
Chrome/Firefox/Edge без polyfill \\
NF5 (детермінованість) & Логічний час + статичний граф & Ідентичні входи
\(\Rightarrow\) ідентичний журнал \\
NF6 (MIT/Apache) & Вибір лише OSS-бібліотек & React, React~Flow,
Zustand, TypeScript \\
NF7 (без встановлення) & SPA у статичному хостингу & Одноразове
завантаження бандлу \\
NF8 (зручність) & React~Flow + панелі Controls/Metrics & Двокліковий
доступ до базових операцій \\
\end{longtable}

\protect\phantomsection\label{_Toc228121450}{}2.2. Математичне та
алгоритмічне забезпечення

У цьому підрозділі подано математичний апарат 16-точкового ШПФ з
основою~4, що становить інтелектуальне ядро симулятора, та формалізацію
обчислень у парадигмі потоку даних.

Дискретне перетворення Фур'є (ДПФ) скінченної комплекснозначної
послідовності \(\{x_{n}\}_{n=0}^{N-1}\) визначається як~(Nussbaumer
1982; Процько et al. 2025):

\[X_{k} = \sum_{n=0}^{N-1}x_{n}\,e^{-j2\pi nk/N} = \sum_{n=0}^{N-1}x_{n}\,W_{N}^{nk},\quad k = 0,1,\ldots,N-1,\]

де \(W_{N}=e^{-j2\pi/N}\) --- поворотний множник (twiddle factor). Пряме
обчислення потребує \(N^{2}\) комплексних множень, тобто має складність
\(O(N^{2})\), що недостатньо для задач реального часу. Клас швидких
алгоритмів (ШПФ) знижує складність до \(O(N\log_{2}N)\) за рахунок
факторизації \(N=N_{1}\cdot N_{2}\). Для \(N=16\) природним вибором є
алгоритм з основою~4 (radix-4), де \(N_{1}=N_{2}=4\)~(Cooley and
Tukey 1965).

Для двовимірної \(4\times4\) декомпозиції часовий і частотний індекси
подаються у вигляді \(n=4s+r\), \(k=4k_{2}+k_{1}\),
де \(r,s,k_{1},k_{2}\in\{0,1,2,3\}\). Вхідна послідовність
\(\{x(n)\}_{n=0}^{15}\) вкладається у матрицю \(\mathbf{A}\) розміром
\(4\times4\): \(A[r][s]=x(4s+r)\). Підстановка у формулу ДПФ із
використанням тотожності \(W_{16}^{4(\cdot)}=W_{4}^{(\cdot)}\) дає
тристадійну схему~(Burrus and Parks 1985):

\[X(4k_{2}+k_{1}) = \sum_{r=0}^{3}W_{16}^{rk_{1}}
\underset{G_{r}(k_{1})}{\underbrace{\Bigl(\sum_{s=0}^{3}A[r][s]\,W_{4}^{sk_{1}}\Bigr)}}
W_{4}^{rk_{2}}.\]

\textbf{Стадія~1 --- рядкові 4-точкові ДПФ.} Для кожного
\(r\in\{0,1,2,3\}\) обчислюється 4-точкове ДПФ рядка:
\(G_{r}(k_{1})=\sum_{s=0}^{3}A[r][s]\,W_{4}^{sk_{1}}\). Чотири ці
перетворення є взаємно незалежними та виконуються паралельно.

\textbf{Стадія~2 --- поворотні множники.} Результати стадії~1
множаться на матрицю поворотних множників:
\(H_{r}(k_{1})=W_{16}^{r\cdot k_{1}}\cdot G_{r}(k_{1})\). При
\(r=0\) або \(k_{1}=0\) маємо \(W_{16}^{0}=1\), тому 7~з 16~множень є
тривіальними. Числові значення 9~нетривіальних множників наведено у
табл.~2.2.

Таблиця 2.2.

Поворотні множники \(W_{16}^{r\cdot k_{1}}\) (округлено до 3~знаків)

\begin{longtable}[]{@{}
  >{\raggedright\arraybackslash}p{(\linewidth - 8\tabcolsep) * \real{0.0978}}
  >{\raggedright\arraybackslash}p{(\linewidth - 8\tabcolsep) * \real{0.0399}}
  >{\raggedright\arraybackslash}p{(\linewidth - 8\tabcolsep) * \real{0.2159}}
  >{\raggedright\arraybackslash}p{(\linewidth - 8\tabcolsep) * \real{0.2384}}
  >{\raggedright\arraybackslash}p{(\linewidth - 8\tabcolsep) * \real{0.2384}}@{}}
\caption[Поворотні множники \(W_{16}^{r\cdot k_{1}}\)]{Поворотні множники
\(W_{16}^{r \cdot k_{1}}\) (округлено до 3 знаків)}\tabularnewline
\toprule\noalign{}
\begin{minipage}[b]{\linewidth}\raggedright
\[r \smallsetminus k_{1}\]
\end{minipage} & \begin{minipage}[b]{\linewidth}\raggedright
0
\end{minipage} & \begin{minipage}[b]{\linewidth}\raggedright
1
\end{minipage} & \begin{minipage}[b]{\linewidth}\raggedright
2
\end{minipage} & \begin{minipage}[b]{\linewidth}\raggedright
3
\end{minipage} \\
\midrule\noalign{}
\endfirsthead
\toprule\noalign{}
\begin{minipage}[b]{\linewidth}\raggedright
\[r \smallsetminus k_{1}\]
\end{minipage} & \begin{minipage}[b]{\linewidth}\raggedright
0
\end{minipage} & \begin{minipage}[b]{\linewidth}\raggedright
1
\end{minipage} & \begin{minipage}[b]{\linewidth}\raggedright
2
\end{minipage} & \begin{minipage}[b]{\linewidth}\raggedright
3
\end{minipage} \\
\midrule\noalign{}
\endhead
\bottomrule\noalign{}
\endlastfoot
0 & \(1\) & \(1\) & \(1\) & \(1\) \\
1 & \(1\) & \(0{,}924 - 0{,}383j\) & \(0{,}707 - 0{,}707j\) &
\(0{,}383 - 0{,}924j\) \\
2 & \(1\) & \(0{,}707 - 0{,}707j\) & \(-j\) & \(-0{,}707 - 0{,}707j\) \\
3 & \(1\) & \(0{,}383 - 0{,}924j\) & \(-0{,}707 - 0{,}707j\) &
\(-0{,}924 + 0{,}383j\) \\
\end{longtable}

\textbf{Стадія~3 --- стовпцеві 4-точкові ДПФ.} Для кожного \(k_{1}\)
обчислюється 4-точкове ДПФ стовпця матриці~\(\mathbf{H}\):
\(X(4k_{1}+k_{2})=\sum_{r=0}^{3}H_{r}(k_{1})\,W_{4}^{rk_{2}}\).
Чотири стовпцевих перетворення також є взаємно незалежними.

Ядром обох стадій є 4-точкове ДПФ, що завдяки властивостям
\(W_{4}\in\{1,-j,-1,j\}\) реалізується без нетривіальних множень. Для
входів \(v_{0},v_{1},v_{2},v_{3}\in\mathbb{C}\) вводяться проміжні
суми і різниці \(E_{0}=v_{0}+v_{2}\), \(E_{1}=v_{0}-v_{2}\),
\(O_{0}=v_{1}+v_{3}\), \(O_{1}=v_{1}-v_{3}\), після чого вихідні
відліки:

\[\begin{aligned}
V_{0} &= E_{0}+O_{0}, & V_{2} &= E_{0}-O_{0}, \\
V_{1} &= E_{1}-j\cdot O_{1}, & V_{3} &= E_{1}+j\cdot O_{1}.
\end{aligned}\]

Множення на \(j\) зводиться до перестановки дійсної та уявної частин зі
зміною знаку (\(j\cdot(a+bj)=-b+aj\)) і не потребує операції множення у
традиційному сенсі. Таким чином, одне 4-точкове ДПФ потребує
\textbf{8~комплексних додавань і 0 нетривіальних множень}~(Nussbaumer 1982).

Загальна обчислювальна вартість алгоритму для \(N=16\): стадія~1 ---
\(4\cdot8=32\) комплексних додавань; стадія~2 --- 9~нетривіальних
комплексних множень (36~дійсних) і 16~комплексних додавань; стадія~3 ---
\(4\cdot8=32\) комплексних додавань. Разом: \textbf{9~нетривіальних
комплексних множень} і \textbf{80~комплексних додавань}, порівняно з
\(N^{2}=256\) множень прямого ДПФ --- прискорення у \(\approx28\)~разів
за кількістю множень~(Nussbaumer 1982; Burrus and Parks 1985).

Алгоритм подається у вигляді орієнтованого ациклічного графа потоку
даних (DFG) \(G=(V,E)\)~(Dennis 1974; Lee and Hurson 1994), де \(V\)
--- множина вузлів чотирьох типів (\texttt{source}, \texttt{dft4},
\texttt{twiddle}, \texttt{sink}), а \(E\subseteq V\times V\) --- множина
спрямованих дуг, що є FIFO-чергами токенів. \textbf{Правило
активації:} вузол~\(v\) є готовим до виконання тоді і лише тоді, коли
на кожному з його вхідних каналів присутній щонайменше один токен; при
активації він поглинає вхідні токени, виконує обчислення та породжує
вихідні токени. \textbf{Часова модель:} вузол має затримку виконання
\texttt{latency}, а дуга --- затримку доставки \texttt{delay}; токен,
породжений у момент~\(t\), стає готовим у момент \(t+\text{delay}\);
логічний час \(\tau\) симулятора інкрементується дискретно на кроці
\(d\tau\), що регулюється множником швидкості.

Планувальник (ядро \texttt{DataflowRuntime}) на кожному кроці симуляції
виконує доставку токенів із FIFO-черг і запуск готових вузлів. Узагальнений
псевдокод наведено нижче.

{\def\LTcaptype{none}
\begin{longtable}[]{@{}
  >{\raggedright\arraybackslash}p{(\linewidth - 0\tabcolsep) * \real{0.5847}}@{}}
\toprule\noalign{}
\endhead
\bottomrule\noalign{}
\endlastfoot
\begin{minipage}[t]{\linewidth}\raggedright
\begin{verbatim}
procedure TICK(dt, maxFires)
    tau <- tau + dt
    for each edge e in E
        while e.queue.head.t + e.delay <= tau
            token <- e.queue.pop()
            v <- target(e); port <- targetPort(e)
            v.inputs[port] <- token
    fires <- 0
    for each node v in V  order by id
        if v is ready and not busy and fires < maxFires
            emit onFire(v, tau)
            v.busy <- tau + v.latency
            outs <- compute(v.kind, v.inputs)
            schedule outs on outgoing edges at time v.busy
            fires <- fires + 1
    for each node v in V
        if v.busy <= tau then v.busy <- 0
\end{verbatim}
\end{minipage} \\
\end{longtable}
}

Псевдокод однієї ітерації ядра симуляції
\texttt{DataflowRuntime.tick()}. Джерело: розроблено автором.

Ця процедура забезпечує детермінованість (NF5): при фіксованому порядку
обходу вузлів результат симуляції не залежить від реального часу
виконання у браузері.

\protect\phantomsection\label{_Toc228121457}{}2.3. Загальна архітектура
системи

Архітектура симулятора описується на двох рівнях моделі C4~(Brown 2018)
--- Context і Container, --- доповнених деталізацією структури DFG.

На контекстному рівні система представляється як «чорна скринька», що
взаємодіє з одним актором --- \emph{користувачем} (студент або
викладач), --- та двома зовнішніми сутностями: \emph{браузером}
(runtime-середовище) і \emph{локальним сховищем браузера} (для
збереження пресетів вхідних сигналів). Застосунок завантажується
одноразово зі статичного веб-хостингу та надалі працює автономно.

На рівні контейнерів симулятор декомпозується на чотири логічні
підсистеми, що виконуються у межах одного процесу браузера:
(1)~\textbf{Simulation Core} (\texttt{DataflowRuntime}, \texttt{scheduler},
\texttt{math}) --- реалізує псевдокод підрозд.~2.2, зберігає стан графа
та черг; (2)~\textbf{Graph Generator} --- модуль побудови DFG для
заданого \(N\) та параметрів затримок; (3)~\textbf{State Store} ---
централізоване сховище стану, до якого підписуються всі UI-компоненти;
(4)~\textbf{UI Layer} --- набір React-компонентів: \texttt{GraphView}
(полотно React~Flow), \texttt{SpectrumView} (діаграма \(|X(k)|\)),
\texttt{Metrics} (пропускна здатність, затримка, черги),
\texttt{Controls} (пуск/пауза/швидкість). UI Layer читає стан зі State
Store та викликає команди; State Store делегує обробку до Simulation
Core; Simulation Core публікує події (\texttt{onToken}, \texttt{onFire},
\texttt{onOutput}), на які State Store оновлює відповідні поля; Graph
Generator викликається один раз при ініціалізації та при кожній зміні
параметрів.

Детальна структура DFG для \(N=16\) складається з \textbf{56~вузлів} і
\textbf{64~дуг}, розподілених за п'ятьма ярусами: 16~вузлів
\texttt{source} (src\(_{0}\ldots\)src\(_{15}\)) генерують токени вхідних
відліків \(x(n)\); 4~вузли \texttt{dft4} (row-0\ldots row-3) виконують
рядкові ДПФ (стадія~1); 16~вузлів \texttt{twiddle} (tw-\(r\)-\(k_{1}\))
множать на поворотні множники (стадія~2); 4~вузли \texttt{dft4}
(col-0\ldots col-3) виконують стовпцеві ДПФ (стадія~3); 16~вузлів
\texttt{sink} (snk\(_{0}\ldots\)snk\(_{15}\)) збирають спектральні
відліки \(X(k)\). Характеристики вузлів і дуг зведено у
табл.~2.3--2.4.

Таблиця 2.3.

Типи вузлів DFG та їхні характеристики

\begin{longtable}[]{@{}
  >{\raggedright\arraybackslash}p{(\linewidth - 10\tabcolsep) * \real{0.1217}}
  >{\raggedright\arraybackslash}p{(\linewidth - 10\tabcolsep) * \real{0.0765}}
  >{\raggedright\arraybackslash}p{(\linewidth - 10\tabcolsep) * \real{0.1270}}
  >{\raggedright\arraybackslash}p{(\linewidth - 10\tabcolsep) * \real{0.1342}}
  >{\raggedright\arraybackslash}p{(\linewidth - 10\tabcolsep) * \real{0.1423}}
  >{\raggedright\arraybackslash}p{(\linewidth - 10\tabcolsep) * \real{0.3983}}@{}}
\caption[Типи вузлів DFG та їхні характеристики]{Типи вузлів DFG та
їхні характеристики}\tabularnewline
\toprule\noalign{}
\begin{minipage}[b]{\linewidth}\raggedright
Тип вузла
\end{minipage} & \begin{minipage}[b]{\linewidth}\raggedright
К-сть
\end{minipage} & \begin{minipage}[b]{\linewidth}\raggedright
Вх. портів
\end{minipage} & \begin{minipage}[b]{\linewidth}\raggedright
Вих. портів
\end{minipage} & \begin{minipage}[b]{\linewidth}\raggedright
Затримка
\end{minipage} & \begin{minipage}[b]{\linewidth}\raggedright
Функція
\end{minipage} \\
\midrule\noalign{}
\endfirsthead
\toprule\noalign{}
\begin{minipage}[b]{\linewidth}\raggedright
Тип вузла
\end{minipage} & \begin{minipage}[b]{\linewidth}\raggedright
К-сть
\end{minipage} & \begin{minipage}[b]{\linewidth}\raggedright
Вх. портів
\end{minipage} & \begin{minipage}[b]{\linewidth}\raggedright
Вих. портів
\end{minipage} & \begin{minipage}[b]{\linewidth}\raggedright
Затримка
\end{minipage} & \begin{minipage}[b]{\linewidth}\raggedright
Функція
\end{minipage} \\
\midrule\noalign{}
\endhead
\bottomrule\noalign{}
\endlastfoot
\texttt{source} & 16 & 0 & 1 & --- & Генерація \(x(n)\) \\
\texttt{dft4} & 8 & 4 & 4 & 400 & 4-точкове ДПФ (метелик radix-4) \\
\texttt{twiddle} & 16 & 1 & 1 & 200 & Множення на \(W_{16}^{rk_{1}}\) \\
\texttt{sink} & 16 & 1 & 0 & --- & Збір \(X(k)\) \\
\multicolumn{2}{@{}>{\raggedright\arraybackslash}p{(\linewidth - 10\tabcolsep) * \real{0.1982} + 2\tabcolsep}}{%
\textbf{Разом: 56 вузлів}} & & & & \\
\end{longtable}

Таблиця 2.4.

Типи дуг DFG

\begin{longtable}[]{@{}
  >{\raggedright\arraybackslash}p{(\linewidth - 6\tabcolsep) * \real{0.1424}}
  >{\raggedright\arraybackslash}p{(\linewidth - 6\tabcolsep) * \real{0.2109}}
  >{\raggedright\arraybackslash}p{(\linewidth - 6\tabcolsep) * \real{0.1435}}
  >{\raggedright\arraybackslash}p{(\linewidth - 6\tabcolsep) * \real{0.1427}}@{}}
\caption[Типи дуг DFG]{Типи дуг DFG}\tabularnewline
\toprule\noalign{}
\begin{minipage}[b]{\linewidth}\raggedright
З вузла
\end{minipage} & \begin{minipage}[b]{\linewidth}\raggedright
До вузла
\end{minipage} & \begin{minipage}[b]{\linewidth}\raggedright
К-сть дуг
\end{minipage} & \begin{minipage}[b]{\linewidth}\raggedright
Затримка
\end{minipage} \\
\midrule\noalign{}
\endfirsthead
\toprule\noalign{}
\begin{minipage}[b]{\linewidth}\raggedright
З вузла
\end{minipage} & \begin{minipage}[b]{\linewidth}\raggedright
До вузла
\end{minipage} & \begin{minipage}[b]{\linewidth}\raggedright
К-сть дуг
\end{minipage} & \begin{minipage}[b]{\linewidth}\raggedright
Затримка
\end{minipage} \\
\midrule\noalign{}
\endhead
\bottomrule\noalign{}
\endlastfoot
\texttt{source} & \texttt{dft4} (рядковий) & 16 & 600 \\
\texttt{dft4} (ряд.) & \texttt{twiddle} & 16 & 600 \\
\texttt{twiddle} & \texttt{dft4} (стовпц.) & 16 & 600 \\
\texttt{dft4} (ст.) & \texttt{sink} & 16 & 600 \\
\multicolumn{2}{@{}>{\raggedright\arraybackslash}p{(\linewidth - 6\tabcolsep) * \real{0.3533} + 2\tabcolsep}}{%
\textbf{Разом: 64 дуги}} & & \\
\end{longtable}

Критичний шлях у DFG (найдовший шлях від \texttt{source} до
\texttt{sink} з урахуванням затримок вузлів і дуг) становить:

\[T_{kr} = 400 + 600 + 200 + 600 + 400 = 2200\]

умовних одиниць часу. Послідовний час виконання (при одному активному
вузлі): \(T_{ps}=8\cdot400+16\cdot200=6400\). Коефіцієнт паралелізму:

\[K_{p} = \frac{T_{ps}}{T_{kr}} = \frac{6400}{2200} \approx 2{,}9,\]

тобто паралельне виконання незалежних вузлів забезпечує майже трикратне
прискорення порівняно з послідовним обчисленням~(Veen 1986).

\protect\phantomsection\label{_Toc228121462}{}2.4. Структури даних

Інформаційне забезпечення симулятора визначається набором типів даних,
що є програмними відповідниками формальних об'єктів DFG.

Усі числові значення подаються структурою \texttt{Complex}
(\(\{re: number,\;im: number\}\)), де \texttt{re} і \texttt{im} ---
дійсна та уявна частини числа \(z=re+j\cdot im\), збережені у подвійній
точності IEEE~754 (64~біти, \(\varepsilon_{mach}\approx2.22\cdot10^{-16}\)).

Тип \texttt{Token} є програмним відповідником токена DFG, що
переміщується по дугам графа. Поля типу наведено у табл.~2.5. Поле
\texttt{t} визначає момент готовності токена: він стає доступним
вузлу-приймачу в момент \(t+\text{edge.delay}\). Поле \texttt{originT}
зберігається незмінним упродовж усього маршруту і використовується для
вимірювання наскрізної затримки.

Таблиця 2.5.

Поля типу \texttt{Token}

\begin{longtable}[]{@{}
  >{\raggedright\arraybackslash}p{(\linewidth - 4\tabcolsep) * \real{0.1004}}
  >{\raggedright\arraybackslash}p{(\linewidth - 4\tabcolsep) * \real{0.1146}}
  >{\raggedright\arraybackslash}p{(\linewidth - 4\tabcolsep) * \real{0.4709}}@{}}
\caption[Поля типу Token]{Поля типу \texttt{Token}}\tabularnewline
\toprule\noalign{}
\begin{minipage}[b]{\linewidth}\raggedright
Поле
\end{minipage} & \begin{minipage}[b]{\linewidth}\raggedright
Тип
\end{minipage} & \begin{minipage}[b]{\linewidth}\raggedright
Призначення
\end{minipage} \\
\midrule\noalign{}
\endfirsthead
\toprule\noalign{}
\begin{minipage}[b]{\linewidth}\raggedright
Поле
\end{minipage} & \begin{minipage}[b]{\linewidth}\raggedright
Тип
\end{minipage} & \begin{minipage}[b]{\linewidth}\raggedright
Призначення
\end{minipage} \\
\midrule\noalign{}
\endhead
\bottomrule\noalign{}
\endlastfoot
\texttt{id} & \texttt{string} & Унікальний ідентифікатор токена \\
\texttt{value} & \texttt{Complex} & Комплексне значення токена \\
\texttt{t} & \texttt{number} & Логічний час надходження на дугу \\
\texttt{originT} & \texttt{number} & Час народження у \texttt{source} \\
\end{longtable}

Тип \texttt{Edge} описує спрямований канал між двома портами вузлів
(табл.~2.6). Поля \texttt{label} і \texttt{labelValue} використовуються
для дуг між стадіями~1 і~2, що несуть мітки поворотних множників
\(W_{16}^{rk_{1}}\).

Таблиця 2.6.

Поля типу \texttt{Edge}

\begin{longtable}[]{@{}
  >{\raggedright\arraybackslash}p{(\linewidth - 4\tabcolsep) * \real{0.1346}}
  >{\raggedright\arraybackslash}p{(\linewidth - 4\tabcolsep) * \real{0.1495}}
  >{\raggedright\arraybackslash}p{(\linewidth - 4\tabcolsep) * \real{0.5600}}@{}}
\caption[Поля типу Edge]{Поля типу \texttt{Edge}}\tabularnewline
\toprule\noalign{}
\begin{minipage}[b]{\linewidth}\raggedright
Поле
\end{minipage} & \begin{minipage}[b]{\linewidth}\raggedright
Тип
\end{minipage} & \begin{minipage}[b]{\linewidth}\raggedright
Призначення
\end{minipage} \\
\midrule\noalign{}
\endfirsthead
\toprule\noalign{}
\begin{minipage}[b]{\linewidth}\raggedright
Поле
\end{minipage} & \begin{minipage}[b]{\linewidth}\raggedright
Тип
\end{minipage} & \begin{minipage}[b]{\linewidth}\raggedright
Призначення
\end{minipage} \\
\midrule\noalign{}
\endhead
\bottomrule\noalign{}
\endlastfoot
\texttt{id} & \texttt{string} & Унікальний ідентифікатор дуги \\
\texttt{from} & \texttt{\{node, port\}} & Вузол і порт-джерело \\
\texttt{to} & \texttt{\{node, port\}} & Вузол і порт-приймач \\
\texttt{delay} & \texttt{number} & Затримка доставки токена (600) \\
\texttt{label} & \texttt{string?} & Мітка на ребрі (\(W_{16}^{rk_{1}}\)) \\
\texttt{labelValue} & \texttt{Complex?} & Числове значення поворотного
множника \\
\end{longtable}

Тип \texttt{NodeSpec} визначає специфікацію вузла DFG, а тип
\(\texttt{Graph}=\{\texttt{nodes: NodeSpec[]},\,\texttt{edges: Edge[]}\}\)
--- повний граф обчислювача. Поля \texttt{NodeSpec} наведено у
табл.~2.7. Для \(N=16\) об'єкт \texttt{Graph} містить 56~специфікацій
вузлів і 64~дуги. Тип \texttt{NodeKind} є переліковим:
\texttt{'source'~|~'dft4'~|~'twiddle'~|~'sink'} і використовується як
ключ у шаблоні Strategy для вибору функції обчислення вузла.

Таблиця 2.7.

Поля типу \texttt{NodeSpec}

\begin{longtable}[]{@{}
  >{\raggedright\arraybackslash}p{(\linewidth - 4\tabcolsep) * \real{0.1117}}
  >{\raggedright\arraybackslash}p{(\linewidth - 4\tabcolsep) * \real{0.1275}}
  >{\raggedright\arraybackslash}p{(\linewidth - 4\tabcolsep) * \real{0.4959}}@{}}
\caption[Поля типу NodeSpec]{Поля типу \texttt{NodeSpec}}\tabularnewline
\toprule\noalign{}
\begin{minipage}[b]{\linewidth}\raggedright
Поле
\end{minipage} & \begin{minipage}[b]{\linewidth}\raggedright
Тип
\end{minipage} & \begin{minipage}[b]{\linewidth}\raggedright
Призначення
\end{minipage} \\
\midrule\noalign{}
\endfirsthead
\toprule\noalign{}
\begin{minipage}[b]{\linewidth}\raggedright
Поле
\end{minipage} & \begin{minipage}[b]{\linewidth}\raggedright
Тип
\end{minipage} & \begin{minipage}[b]{\linewidth}\raggedright
Призначення
\end{minipage} \\
\midrule\noalign{}
\endhead
\bottomrule\noalign{}
\endlastfoot
\texttt{id} & \texttt{string} & Ідентифікатор вузла \\
\texttt{kind} & \texttt{NodeKind} & Тип (source, dft4, twiddle, sink) \\
\texttt{inPorts} & \texttt{string[]} & Список вхідних портів \\
\texttt{outPorts} & \texttt{string[]} & Список вихідних портів \\
\texttt{latency} & \texttt{number} & Тривалість обчислення (400 або 200) \\
\texttt{params} & \texttt{object?} & Параметри (для twiddle: \(N,k\)) \\
\end{longtable}

Алгоритм оперує чотирма матрицями даних розміром \(4\times4\):

\[\begin{aligned}
\mathbf{A}[r][s] &= x(4s+r), & &\text{вхідні дані,} \\
\mathbf{G}[r][k_{1}] &= DFT4(\mathbf{A}[r][\cdot])[k_{1}], & &\text{після стадії 1,} \\
\mathbf{H}[r][k_{1}] &= W_{16}^{r\cdot k_{1}}\cdot\mathbf{G}[r][k_{1}], & &\text{після стадії 2,} \\
X(4k_{1}+k_{2}) &= DFT4(\mathbf{H}[\cdot][k_{1}])[k_{2}]. & &\text{вихідний спектр.}
\end{aligned}\]

Для контрольної верифікації використано лінійний вхідний сигнал
\(x(n)=n\), \(n=0\ldots15\). Проміжні значення для першого рядка
(\(r=0\)) і перших чотирьох спектральних відліків наведено у
табл.~2.8. Точне значення \(X(0)=\sum_{n=0}^{15}n=120\) підтверджує
коректність стадії~3; рядок \(\mathbf{H}[0][\cdot]\) збігається з
\(\mathbf{G}[0][\cdot]\), оскільки \(W_{16}^{0\cdot k_{1}}=1\) для
всіх~\(k_{1}\).

Таблиця 2.8.

Проміжні значення алгоритму для \(x(n)=n\), рядок \(r=0\)

\begin{longtable}[]{@{}
  >{\raggedright\arraybackslash}p{(\linewidth - 8\tabcolsep) * \real{0.2233}}
  >{\raggedright\arraybackslash}p{(\linewidth - 8\tabcolsep) * \real{0.1370}}
  >{\raggedright\arraybackslash}p{(\linewidth - 8\tabcolsep) * \real{0.1733}}
  >{\raggedright\arraybackslash}p{(\linewidth - 8\tabcolsep) * \real{0.1733}}
  >{\raggedright\arraybackslash}p{(\linewidth - 8\tabcolsep) * \real{0.1733}}@{}}
\caption[Проміжні значення алгоритму]{Проміжні значення алгоритму для
\(x(n)=n\), рядок \(r=0\)}\tabularnewline
\toprule\noalign{}
\begin{minipage}[b]{\linewidth}\raggedright
Матриця
\end{minipage} & \begin{minipage}[b]{\linewidth}\raggedright
Індекс 0
\end{minipage} & \begin{minipage}[b]{\linewidth}\raggedright
Індекс 1
\end{minipage} & \begin{minipage}[b]{\linewidth}\raggedright
Індекс 2
\end{minipage} & \begin{minipage}[b]{\linewidth}\raggedright
Індекс 3
\end{minipage} \\
\midrule\noalign{}
\endfirsthead
\toprule\noalign{}
\begin{minipage}[b]{\linewidth}\raggedright
Матриця
\end{minipage} & \begin{minipage}[b]{\linewidth}\raggedright
Індекс 0
\end{minipage} & \begin{minipage}[b]{\linewidth}\raggedright
Індекс 1
\end{minipage} & \begin{minipage}[b]{\linewidth}\raggedright
Індекс 2
\end{minipage} & \begin{minipage}[b]{\linewidth}\raggedright
Індекс 3
\end{minipage} \\
\midrule\noalign{}
\endhead
\bottomrule\noalign{}
\endlastfoot
\(\mathbf{A}[0][\cdot]\) & \(0\) & \(4\) & \(8\) & \(12\) \\
\(\mathbf{G}[0][\cdot]\) & \(24+0j\) & \(-8+8j\) & \(-8+0j\) & \(-8-8j\) \\
\(\mathbf{H}[0][\cdot]\) & \(24+0j\) & \(-8+8j\) & \(-8+0j\) & \(-8-8j\) \\
\(X(k)\), \(k=0\ldots3\) & \(120+0j\) & \(\approx-8+40j\) &
\(\approx-8+19j\) & \(\approx-8+12j\) \\
\end{longtable}

\protect\phantomsection\label{_Toc228121469}{}2.5. Механізми обміну
даними та інтеграції

Оскільки симулятор не має серверної частини, класичні механізми мережевої
інтеграції (REST API, WebSocket, gRPC) у ньому не застосовуються.
Координація підсистем реалізується через \emph{внутрішній подійний
канал} і \emph{реактивне сховище стану}.

Ядро \texttt{DataflowRuntime} публікує три типи подій, що утворюють
внутрішній API симулятора (табл.~2.9). Події реалізують шаблон
Publisher/Subscriber: видавець (ядро) не має прямих посилань на
підписників, а лише передає об'єкт події у список зареєстрованих
обробників. Це забезпечує ізоляцію бізнес-логіки від UI та дозволяє
додавати нові підписники (наприклад, експорт event~log у файл) без
модифікації ядра.

Таблиця 2.9.

Події шини симуляції

\begin{longtable}[]{@{}
  >{\raggedright\arraybackslash}p{(\linewidth - 4\tabcolsep) * \real{0.1201}}
  >{\raggedright\arraybackslash}p{(\linewidth - 4\tabcolsep) * \real{0.3574}}
  >{\raggedright\arraybackslash}p{(\linewidth - 4\tabcolsep) * \real{0.5226}}@{}}
\caption[Події шини симуляції]{Події шини симуляції}\tabularnewline
\toprule\noalign{}
\begin{minipage}[b]{\linewidth}\raggedright
Подія
\end{minipage} & \begin{minipage}[b]{\linewidth}\raggedright
Умова виникнення
\end{minipage} & \begin{minipage}[b]{\linewidth}\raggedright
Підписники
\end{minipage} \\
\midrule\noalign{}
\endfirsthead
\toprule\noalign{}
\begin{minipage}[b]{\linewidth}\raggedright
Подія
\end{minipage} & \begin{minipage}[b]{\linewidth}\raggedright
Умова виникнення
\end{minipage} & \begin{minipage}[b]{\linewidth}\raggedright
Підписники
\end{minipage} \\
\midrule\noalign{}
\endhead
\bottomrule\noalign{}
\endlastfoot
\texttt{onToken} & Токен надійшов на дугу \(e\in E\) & Heatmap черг,
Metrics (queue size) \\
\texttt{onFire} & Активовано вузол \(v\in V\) & Heatmap активності
вузлів \\
\texttt{onOutput} & Токен надійшов до \texttt{sink} & SpectrumView,
Metrics (throughput, latency) \\
\end{longtable}

Стан симулятора централізований у сховищі \texttt{simStore}, що містить
чотири групи полів: \emph{стан графа} (\texttt{graph: Graph},
\texttt{queues: Map<EdgeId, Token[]>},
\texttt{nodeBusy: Map<NodeId, number>}); \emph{стан виконання}
(\texttt{status: 'IDLE'|'RUNNING'|'PAUSED'|'DONE'}, \texttt{simTime},
\texttt{speed}); \emph{метрики} (\texttt{throughputEMA},
\texttt{latencyEMA}, \texttt{queueSizeEMA},
\texttt{activity: Map<NodeId, number>}); \emph{результат}
(\texttt{spectrum: Complex[16]}, \texttt{eventLog: Event[]}). UI-компоненти
підписуються на окремі поля сховища за допомогою селекторів і
перемальовуються лише при зміні підписаних полів, що забезпечує стабільні
60~fps (NF3) навіть при інтенсивному потоці подій.

Для підтримки повторюваних демонстраційних сценаріїв (UC1) передбачено
механізм пресетів. Пресет --- це JSON-об'єкт, що містить вхідний
сигнал \(\{x_{n}\}\), значення \texttt{latency} і \texttt{delay} та
початкову швидкість симуляції. Пресети зберігаються у Local~Storage
браузера за ключем \texttt{simulator.presets} і завантажуються через
компонент \texttt{PresetSelector}. Експорт пресету і журналу подій (UC9)
реалізується стандартним Web~API \texttt{Blob+URL.createObjectURL} без
передачі на сервер, що відповідає принципу конфіденційності
(підрозд.~1.4).

\protect\phantomsection\label{_Toc228121473}{}2.6. Інтерфейс
користувача та взаємодія із системою

Інтерфейс симулятора спроєктовано за принципом \emph{інформаційної
ієрархії}: найбільша площа екрана відводиться під центральний об'єкт
дослідження --- граф потоку даних, --- а допоміжні панелі (метрики,
спектр, керування) розміщуються по периметру, не перекриваючи граф.
Головний екран розділено на чотири функціональні зони.

\textbf{Центральна зона (GraphView)} --- полотно React~Flow з вузлами та
анімованими дугами; користувач може перетягувати вузли, змінювати масштаб
колесом миші та клікати по вузлу для перегляду поточних значень на
входах/виходах. \textbf{Верхня панель (Controls)} --- кнопки «Пуск»,
«Пауза», «Крок», «Скидання», повзунок множника швидкості (\(0.1\times\)
до \(10\times\)), селектор пресету вхідного сигналу і поле редагування
\texttt{latency}/\texttt{delay}. \textbf{Права панель (SpectrumView +
Metrics)} --- стовпчикова діаграма \(|X(k)|\) у лінійному або
логарифмічному масштабі та три числових індикатори (throughput, latency,
queue size). \textbf{Нижня панель (TerminalOutput)} --- журнал подій у
реальному часі з фільтрацією за типом та можливістю експорту у файл.

Базовий user-flow для студента (UC3) складається з трьох кроків,
доступних у межах двох кліків від головного екрана (NF8):
(1)~вибір пресету вхідного сигналу або введення власних відліків через
\texttt{InputEditor};
(2)~натискання кнопки «Пуск» --- ядро починає викликати \texttt{tick()}
у циклі \texttt{requestAnimationFrame} з кроком \(d\tau=16.67\)~мс
\(\cdot\)~speed;
(3)~спостереження за рухом токенів у GraphView, зміною метрик та
поступовим заповненням SpectrumView. У режимі паузи доступний покроковий
режим~\emph{Step}: одне натискання ініціює один виклик \texttt{tick()},
що дозволяє детально дослідити виконання алгоритму.

Для підсилення сприйняття потокової природи алгоритму UI використовує
три допоміжні візуальні механізми. \textbf{Анімація токенів:} кожен
токен відображається як кружечок, що переміщується по дузі в інтервалі
\([t,\,t+\text{delay}]\) з плавною інтерполяцією позиції; числове
значення (дійсна та уявна частини) виводиться безпосередньо на кружечку.
\textbf{Теплова карта активності вузлів:} заливка кожного вузла
модулюється значенням \(a_{i}(\tau)\), що оновлюється за
експоненціальним законом
\(a_{i}(\tau+d\tau)=\alpha\cdot a_{i}(\tau)+\Delta a_{i}\), де
\(\alpha=0{,}92\) і \(\Delta a_{i}=1\) при подіях \texttt{onFire}; вузли,
що нещодавно часто активувалися, відображаються яскравішими.
\textbf{Стан busy на вузлах:} поки вузол виконує обчислення
(\(v.busy>\tau\)), навколо нього відображається кільцевий індикатор
прогресу, що ілюструє затримку \texttt{latency}.

Загальний вигляд інтерфейсу подано на рис.~2.1.

\protect\phantomsection\label{fig:ui}{}Рис.~2.1. Макет інтерфейсу
веб-симулятора потокового 16-точкового ШПФ: центральна зона GraphView,
верхня панель Controls, права панель SpectrumView + Metrics, нижня панель
TerminalOutput. Джерело: розроблено автором.

\protect\phantomsection\label{_Toc228121478}{}2.7. Висновки до розділу~2

У другому розділі спроєктовано програмне рішення, що відповідає вимогам,
сформульованим у підрозд.~1.5.

Обрано архітектурний підхід клієнтського SPA без серверної частини, що
одночасно забезпечує виконання вимог NF2 (відгук \(\leq100\)~мс), NF4
(портативність Chrome/Firefox/Edge), NF7 (робота без встановлення) і NF5
(детермінованість). Вибір реактивного рендерингу та централізованого
сховища стану обґрунтований як компроміс між продуктивністю та
супроводжуваністю коду. Застосовано шаблони Observer, State, Strategy,
Factory і Facade; зв'язок між вимогами та рішеннями показано у
табл.~2.1.

Побудовано математичну модель алгоритму ШПФ radix-4 для \(N=16\) у
вигляді тристадійної \(4\times4\) декомпозиції, що потребує
9~нетривіальних комплексних множень і 80~комплексних додавань проти
256~множень прямого ДПФ. Сформульовано формальну модель потоку даних із
правилом активації та часовою семантикою, подано псевдокод планувальника.

Спроєктовано архітектуру системи на рівнях C4 Context і Container з
декомпозицією на чотири підсистеми (Simulation Core, Graph Generator,
State Store, UI Layer). Деталізовано граф обчислювача: 56~вузлів
чотирьох типів і 64~дуги з затримкою 600~одиниць, критичним шляхом
\(T_{kr}=2200\) і коефіцієнтом паралелізму \(K_{p}\approx2{,}9\).

Визначено структури даних \texttt{Token}, \texttt{Edge}, \texttt{NodeSpec},
\texttt{Graph} --- програмні відповідники формальних об'єктів DFG;
комплексні числа подаються у форматі IEEE~754 double, що гарантує запас
точності відносно \(\varepsilon=10^{-9}\) (NF1). Обрано механізми обміну
даними: внутрішня шина подій (\texttt{onToken}, \texttt{onFire},
\texttt{onOutput}) для ізоляції UI від ядра, реактивне сховище та
пресети у Local~Storage для повторюваних сценаріїв. Спроєктовано
інтерфейс з інформаційною ієрархією, трьома візуальними механізмами
відображення потокової природи алгоритму та базовим user-flow у межах
двох кліків (NF8).

Описані проєктні рішення утворюють повне технічне завдання для
реалізації, що подається у розділі~3.